\begin{center}
    \textbf{ABSTRACT}
\end{center}

[cite_start]Partial Discharge (PD) is a localized dielectric breakdown in high-voltage insulation systems and a key indicator of potential equipment failure[cite: 12]. [cite_start]The reliability of power systems hinges on the early and accurate detection of PD, but traditional manual analysis is subjective and inefficient, while automated methods face challenges with feature engineering, noise, and model generalizability[cite: 13].

[cite_start]This research addresses a critical gap in the existing literature: the lack of systematic, comparative analysis of traditional machine learning (ML) algorithms across diverse PD datasets and feature sets[cite: 14]. [cite_start]This study conducts a comprehensive comparison of five supervised ML algorithms---Support Vector Classifier (SVC), Random Forest (RF), K-Nearest Neighbors (KNN), Decision Tree (DT), and an Artificial Neural Network (ANN)---to diagnose insulation defects[cite: 15].

[cite_start]A multi-domain feature extraction strategy is employed, deriving features from the time, statistical, frequency, and time-frequency (using Discrete Wavelet Transform) domains to create a robust representation of PD events[cite: 16]. [cite_start]The performance of these classifiers is rigorously evaluated using a 10-fold cross-validation strategy on a long-term dataset of PD signals from overhead power lines[cite: 17]. [cite_start]Key performance metrics, including accuracy, precision, recall, F1-score, and computational efficiency, are used to identify optimal feature-algorithm pairings and assess model generalizability[cite: 18].

[cite_start]The aim is to provide a foundational benchmark for selecting robust and efficient diagnostic solutions for real-world field conditions, overcoming the limitations of studies based on single, homogeneous datasets[cite: 19].